%% =====================================================================
%% DRAFT replacement for SI Section S3 "MCAP Coverage Study"
%% (si.Rnw lines ~282-368). NOT yet inserted into si.Rnw — for Bo's review.
%% Produced 2026-06-03. Drop-in: keeps \label{sec:mcap-coverage},
%% Table:mcap-coverage, Table:three-loglik, fig:mcap-coverage-plots, and the
%% \native/\invasive notation. Wrapped in \bye{} per the revision convention.
%%
%% OPEN ITEMS FLAGGED INLINE BELOW (search "VERIFY") — resolve before insertion:
%%   1. ri/probi/f_Si "product ridge" algebra vs the actual C-snippet.
%%   2. sigma_F undercovers with ~zero MLE drift -> width, not displacement.
%%   3. within-chain-SD <-> coverage link is not monotone (sigP vs sigF).
%% =====================================================================

\bye{We assess the empirical coverage of Monte Carlo Adjusted Profile (MCAP) confidence intervals \citep{ionides17} on the SIRJPF2 model by parametric bootstrap, leaning on simulation rather than on the asymptotic local-normality approximation that nominally justifies MCAP.
We simulated $B=100$ datasets from the all-shared SIRJPF2 maximum likelihood estimate and, on each simulated dataset, recomputed the MCAP profile and its nominal $95\%$ confidence interval for four target parameters chosen to span the identifiability range: the native and invasive birth rates $r^{\native}$ and $r^{\invasive}$, the parasite spore process noise $\sigma_P$, and the food process noise $\sigma_F$.
Profiles were computed with the same iterated-filtering procedure used for the data analysis in the main paper, at the production run level.
A confidence interval ``covers'' on a given simulated dataset if it contains the known data-generating value of the target parameter, and empirical coverage is the fraction of the $B$ datasets on which it does.
We report this coverage honestly, with its Monte Carlo standard error, and we use the bootstrap to diagnose, rather than to conceal, the gap between empirical and nominal coverage.}

\subsection{Empirical coverage}

\bye{Table~\ref{Table:mcap-coverage} summarises the headline result and Figure~\ref{fig:mcap-coverage-plots} displays the per-dataset intervals.
For the three identifiable target parameters the nominal $95\%$ MCAP interval covers the truth in a clear majority of simulations, but below the nominal level: $77.8\%$ for $r^{\native}$ ($77/99$; one profile failed to converge and is excluded, so the denominator is $99$), $87.0\%$ for $\sigma_P$ ($87/100$), and $74.0\%$ for $\sigma_F$ ($74/100$).
With these denominators the Monte Carlo standard error of a coverage estimate, $\sqrt{\hat c(1-\hat c)/n}$, is about $4.2$ percentage points for $r^{\native}$, $3.4$ for $\sigma_P$, and $4.4$ for $\sigma_F$.
The shortfalls for $r^{\native}$ and $\sigma_F$ therefore sit roughly four to five standard errors below $0.95$, and that for $\sigma_P$ about two and a half; the departures for $r^{\native}$ and $\sigma_F$ are real sub-nominal coverage, not bootstrap noise.

%% VERIFY (2): The coverage-table mean MCAP estimate is biased HIGH for r^n and
%% LOW for sigma_P/sigma_F. But the 5-dataset diagnostic drift for sigma_F is
%% only -0.10 (near zero) -- i.e. sigma_F's MLE is barely displaced yet it
%% still undercovers, which points to interval WIDTH, not MLE displacement, for
%% that parameter. The single "MLE displacement" story below fits r^n cleanly
%% but NOT sigma_F. Either (a) soften to "displacement and/or finite-sample
%% interval miscalibration", or (b) lead with displacement for r^n and note
%% sigma_F separately as a width effect that cutoff calibration (option B) fixes.
For $r^{\native}$ the mean MCAP point estimate is displaced from the truth in a consistent direction (biased high; Table~\ref{Table:mcap-coverage}), indicating that its undercoverage is driven by displacement of the maximum likelihood estimate rather than by intervals that are too narrow; for $\sigma_F$ the point estimate is nearly unbiased, so its shortfall instead reflects intervals that are slightly too narrow in finite samples --- a distinction we return to below.

For $r^{\invasive}$ the situation is qualitatively different and must not be read as well-calibrated coverage.
%% VERIFY (1): Algebra of the ri/probi/f_Si ridge. The model C-snippet has
%%   Ji_term  = ri    * f_Si * F * Si          (recruitment: ri * f_Si)
%%   Ii_term  = probi * f_Si * Si * P          (infection:  probi * f_Si)
%% so f_Si is SHARED between two SEPARATE pairwise products (ri*f_Si and
%% probi*f_Si), and f_Si ALSO appears alone in the F and P depletion terms.
%% The single triple product "r^i f_S^i p^i" below is therefore imprecise.
%% Suggested correct phrasing: "r^i is confounded with f_S^i through the
%% recruitment product r^i f_S^i, and p^i with f_S^i through the infection
%% product p^i f_S^i; f_S^i is itself only weakly pinned (it enters alone only
%% in the food/parasite depletion terms), leaving all three weakly identified."
%% CONFIRM against the model before this goes to Ed/referees.
In the SIRJPF2 process model $r^{\invasive}$, $f_S^{\invasive}$ and $p^{\invasive}$ are jointly weakly identified through the invasive-species recruitment and infection products, so the level of $r^{\invasive}$ is individually non-identifiable.
The corresponding profile is essentially flat over the parameter grid (Figure~\ref{fig:mcap-coverage-plots}d), each interval spans the grid, and the nominal $95.8\%$ ``coverage'' is a degenerate artifact of an uninformative interval rather than evidence of correct calibration.
We therefore present $r^{\invasive}$ only to document this ridge and base our coverage conclusions on the three identifiable parameters.
The contrast is itself the diagnostic: it localises the residual coverage gap to weakly-identified directions of the likelihood surface rather than to MCAP as a method.}

\subsection{Undercoverage is a finite-sample, not a numerical, phenomenon}

\bye{The sub-nominal coverage of $r^{\native}$, $\sigma_P$ and $\sigma_F$ is at first glance concerning, since profile-likelihood intervals are asymptotically justified by local asymptotic normality.
Two competing explanations must be separated before any correction is contemplated: either the profiles are under-converged, a numerical Monte Carlo artifact that additional computation could remove, or the maximum likelihood estimate is genuinely displaced from the truth in finite samples, a property of the model and sample size that no amount of further filtering will change.
We discriminate between these with a three-log-likelihood diagnostic on five representative simulated datasets ($b \in \{2, 25, 50, 75, 95\}$).
On each dataset we evaluated $\ell_{\text{truth}}$, the particle-filter log-likelihood at the data-generating $\theta_{\text{true}}$ with no further optimisation; $\ell_{\text{prof}}^{\max}$, the best log-likelihood across all points of the production coverage profile; and $\ell_{\text{glob}}^{\text{best}}$, the best log-likelihood from a $20$-chain iterated-filtering search with all parameters free, warm-started from $\theta_{\text{true}}$.

On every dataset the estimated maximum likelihood lies above the truth (Table~\ref{Table:three-loglik}).
A fully free re-optimisation warm-started at the truth clears $\ell_{\text{truth}}$ by $\overline{\ell_{\text{glob}}^{\text{best}} - \ell_{\text{truth}}} \approx 4.5$ log units, and the more thorough production profiles clear it by $\overline{\ell_{\text{prof}}^{\max} - \ell_{\text{truth}}} \approx 8.5$.
These are two estimates of the same quantity, the finite-sample displacement $\ell_{\text{MLE}} - \ell_{\text{truth}}$; the profile is the more thorough optimiser, which is why it exceeds the $20$-chain search ($\overline{\ell_{\text{glob}}^{\text{best}} - \ell_{\text{prof}}^{\max}} \approx -4$, a negative gap).
That negative gap is the key point: the production profiles already search the surface at least as well as a free re-optimisation started at the answer, so they are not leaving log-likelihood on the table and additional particles or chains would not raise them --- the numerical explanation is ruled out.
Both estimates lie below the asymptotic Wilks expectation $\mathbb{E}[\ell_{\text{MLE}} - \ell_{\text{truth}}] = p/2 = 12$, where SIRJPF2 has $p = 24$ free interior parameters (the parameter vector has $26$ entries, with $\sigma_{S^{\native}}$ and $\sigma_{S^{\invasive}}$ fixed at the boundary value zero) fit to an $80$-observation panel ($U = 8$ units, $N_u = 10$ time points each); a displacement of this order, below the full asymptotic amount, is exactly what one expects for a weakly-identified model on a short panel, where along flat directions the data carry little information.
(Both gaps are maxima of Monte Carlo log-likelihood estimates and so carry a modest upward selection bias; MCAP accounts for this Monte Carlo uncertainty when it forms the interval, so it does not distort the coverage construction itself.)
The undercoverage is therefore a genuine finite-sample property of the estimator along weakly-identified ridges rather than a failure of the MCAP construction: an MCAP interval centred near each dataset's displaced MLE will, on these datasets, sometimes miss a truth that sits below it on the ridge.}

\subsection{Identifiability findings}

\bye{Consistent with this diagnosis, the largest finite-sample effects are concentrated in the weakly-identified parameters.
Across the $20$ iterated-filtering chains used in the diagnostic we examined the within-chain dispersion of all $24$ free parameters, which gives a direct, model-internal map of identifiability.
The ridge parameters $r^{\invasive}$, $f_S^{\invasive}$ and $p^{\invasive}$ each show within-chain standard deviations of $\approx 2.5$ log units, reflecting the joint non-identifiability noted above; this is the same ridge responsible for the degenerate $r^{\invasive}$ coverage in panel~(d) of Figure~\ref{fig:mcap-coverage-plots}.
Two further parameters, $\sigma_{I^{\invasive}}$ and $\sigma_{I^{\native}}$, are dispersed by $\approx 5$ log units each (a factor of $\approx 120$) and we treat them as essentially non-identifiable.
By contrast the better-identified parameters $r^{\native}$, $f_S^{\native}$, $p^{\native}$ and $\sigma_F$ have within-chain standard deviations of only $1.4$--$1.5$ log units, with finite-sample log-MLE drift of order $\pm 1$; $\sigma_P$ is intermediate ($\approx 2.3$).
%% VERIFY (3): within-chain SD does NOT line up monotonically with coverage:
%% sigma_F has the SMALLEST within-SD (1.44) but the WORST coverage (74%), while
%% sigma_P has a LARGER within-SD (2.29) but BETTER coverage (87%). So do NOT
%% claim "flattest surface => lowest coverage" as a clean law. Keep within-SD as
%% an identifiability MAP, and attribute the rn shortfall to MLE displacement and
%% the sigma_F shortfall to interval width (see VERIFY (2)), rather than forcing
%% a single monotone story.
The within-chain dispersion is thus an identifiability map rather than a direct predictor of coverage: the ridge directions are where the estimator is least constrained, while the precise coverage of each identifiable parameter reflects the interaction of its (small) finite-sample bias with the width of its MCAP interval.}

\subsection{Interpretation and reporting}

\bye{We read these results as an honest characterisation of inferential uncertainty rather than as an invalidation of the analysis.
The MCAP intervals remain asymptotically justified and numerically well-converged; the sub-nominal coverage is the finite-sample cost of fitting a richly parameterised mechanistic model, with several weakly-identified directions, to a short ecological panel.
The SIRJPF2 analysis is in this sense exploratory rather than confirmatory, and we mean that operationally: the study is designed to characterise plausible mechanisms and to map which parameters the panel identifies and which it does not, not to deliver certified $95\%$-coverage guarantees for individual weakly-identified parameters for downstream decision or regulatory use.
For the well-identified parameters the bootstrap shows MCAP coverage close to nominal, and for these the nominal intervals can be read at face value.
For the weakly-identified parameters the residual undercoverage is one further, quantified component of the uncertainty already acknowledged for this model, to be weighed alongside structural and model uncertainty rather than hidden.
Because the diagnostic establishes that the shortfall is a reproducible finite-sample property of the estimator and not a numerical artifact, the same parametric-bootstrap distribution that exposed it can also be used to widen the MCAP cutoff for an identifiable parameter to the empirical quantile that restores $95\%$ coverage; this is a single calibration of an existing estimator against its own parametric-bootstrap distribution, not a nested bootstrap-of-bootstraps.
We caution, however, that calibrating a $95\%$ quantile from $B=100$ datasets is itself noisy, so we report the uncalibrated empirical coverage with its Monte Carlo standard error as the primary, fully transparent result, and treat any cutoff calibration as a sensitivity check rather than as a replacement for the reported intervals.}
